
\section{Introduction}
\label{sec:introduction}

Tensor renormalization group (TRG) is a promising computational approach to study the lattice field theory. 
The dynamics of the theory can be investigated by the TRG
mostly in the thermodynamic limit without suffering from the negative sign problem
since it does not employ any stochastic process. 
The TRG was originally proposed by Levin and Nave as a 
real space renormalization group 
for the two-dimensional Ising model \cite{Levin:2006jai}. 
Extensions to fermionic systems were firstly discussed by Gu {\it et al.} \cite{Gu:2010yh,Gu:2013gba},  
and the Grassmann TRG has been applied to many models such as
the Schwinger model with and without $\theta$ term 
\cite{Shimizu:2014uva,Shimizu:2014fsa,Shimizu:2017onf}, 
the Gross-Neveu model with finite density \cite{Takeda:2014vwa}, 
and $\mathcal{N}=1$ Wess-Zumino model \cite{Kadoh:2018hqq}.
\footnote{See Refs.~\cite{Sakai:2017jwp,Yoshimura:2017jpk,Meurice:2018fky,Butt:2019uul} for other related studies. }
These earlier works have verified that the TRG is also useful to evaluate 
the path integral over the Grassmann variables. 
 
In this paper, we introduce a different way from these studies to derive a tensor network 
representation and to implement the Grassmann TRG. The tensor network 
is derived with auxiliary Grassmann fields. 
As an advantage over the conventional tensor network formulation, 
this method allows us to find a general formula, characterized by the singular 
value decomposition for the Dirac matrix, to obtain the tensor network representation for any lattice 
fermion theory with only nearest neighbor interactions. 


This paper is organized as follows. In Sec. \ref{sec:formulation}, we define a 
Grassmann tensor and its contraction rule. We explain how to construct the 
tensor network with auxiliary Grassmann fields in Sec. \ref{sec:TN}. 
To see the validity of the current formulation, numerical results for the two-dimensional free 
Wilson fermion are provided in Sec. \ref{sec:GTRG}. Sec. \ref{sec:summary} is 
devoted to summary and outlook.

